% 扫描《大学物理简明教程》第5版第12章量子物理基础
% 选择12.1.1—12.1.14，原书/PDF316—317页。原创推导，原题题意简述。
\subsection{选择题：量子辐射、光子、氢谱与量子态}

\sourceproblem{12.1.1}{316}{316}
\problemtext 绝对黑体应具备怎样的吸收/反射性质？
\workings
绝对黑体按定义吸收入射的\emph{全部波长}电磁辐射，
无反射、无透射，故题中“不能反射任何光线”正确，选\textbf{D}。
\crosscheck “黑体”不等于在任意温度下不辐射可见光；
热黑体仍会按温度发射热辐射。

\sourceproblem{12.1.2}{316}{316}
\problemtext 钠的光电红限波长$540$ nm，某单色光使最大光电子动能为$1.2$ eV；
求入射光的波长。
\workings
逸出功$W=hc/(540\,\mathrm{nm})$，
爱因斯坦方程$hc/\lambda=W+1.2$ eV，
采用$hc\simeq1240\,\mathrm{eV\,nm}$，
\[
\boxed{\lambda=\frac{1240}{1240/540+1.2}\,\mathrm{nm}
\simeq354.6\,\mathrm{nm}}.
\]
选\textbf{D（355 nm）}。
\crosscheck 入射光要比红限更短波，结果满足$\lambda<540$ nm。

\sourceproblem{12.1.3}{316}{316}
\problemtext 将地球与金星均视为理想黑体，温度分别$293$ K和$773$ K，
且两者尺寸近似相同，比较单位表面面积的辐射出射度。
\workings
斯特藩--玻尔兹曼律$M=\sigma T^4$，故
\[
\boxed{\frac{M_{\rm Venus}}{M_{\rm Earth}}
=\left(\frac{773}{293}\right)^4\simeq48.44}.
\]
选\textbf{D（48.4）}。
\crosscheck 两者\emph{单位面积}比值与球体半径无关；对总辐射功率则还需面积。

\sourceproblem{12.1.4}{316}{316}
\problemtext 同种金属的红限频率$\nu_0$，以频率$\nu_1,\nu_2$照射，
第二次遏止电压是第一次的两倍；求$\nu_2$与$\nu_1,\nu_0$的关系。
\workings
$eU_a=h(\nu-\nu_0)$，令$U_{a,2}=2U_{a,1}$，
得$\nu_2-\nu_0=2(\nu_1-\nu_0)$，故
\[
\boxed{\nu_2=2\nu_1-\nu_0}.
\]
选\textbf{C}。
\crosscheck 这不是频率直接加倍，因逸出功引入频率阈值。

\sourceproblem{12.1.5}{316}{316}
\problemtext 康普顿散射中散射后波长为入射波长的1.2倍；
求散射光子能量与反冲电子动能之比。
\workings
光子能量反比波长，$E'_{\gamma}=E_\gamma/1.2=5E_\gamma/6$；
电子原先静止时反冲动能$K=E_\gamma-E'_\gamma=E_\gamma/6$。
\[
\boxed{E'_\gamma/K=5}.
\]
选\textbf{D}。
\crosscheck 忽略外界束缚能；能量守恒适用于自由电子的康普顿散射。

\sourceproblem{12.1.6}{316}{316}
\problemtext 氢原子巴耳末线系中波长最长与次长的谱线分别记$\lambda_1,\lambda_2$，
求两者比值。
\workings
最长线$n=3\to2$，次长线$n=4\to2$，
由$1/\lambda=R_H(1/2^2-1/n^2)$，
\[
\boxed{\frac{\lambda_1}{\lambda_2}
=\frac{1/4-1/16}{1/4-1/9}=\frac{27}{20}}.
\]
选\textbf{C}。
\crosscheck 最大波长应对应能级差最小的巴耳末跃迁。

\sourceproblem{12.1.7}{317}{317}
\problemtext 静止氢原子由$n=3$直接跃迁到$n=1$，
原子质量$1.67\times10^{-27}$kg，估算反冲速度。
\workings
光子能量$\Delta E=13.6(1-1/9)=12.09$ eV，
动量$p_\gamma\simeq\Delta E/c$；
原子反冲动量大小相等，
\[
\boxed{v_r\simeq\frac{12.09(1.602\times10^{-19})}
{(1.67\times10^{-27})(3.00\times10^8)}
\simeq3.87\,\mathrm{m/s}}.
\]
选\textbf{A（4 m/s）}。
\crosscheck 原子反冲能远小于跃迁光子能量，忽略它只产生极小修正。

\sourceproblem{12.1.8}{317}{317}
\problemtext 氢原子由能级$E_3$直接到$E_1$发出$\lambda_3$，
分两步$E_3\to E_2\to E_1$发出$\lambda_2,\lambda_1$；
求波长关系。
\workings
能量守恒
$hc/\lambda_3=hc/\lambda_2+hc/\lambda_1$，
所以
\[
\boxed{\lambda_3^{-1}=\lambda_1^{-1}+\lambda_2^{-1}}.
\]
选\textbf{C}。
\crosscheck 光子\emph{频率/波数}可相加，波长不可直接相加。

\sourceproblem{12.1.9}{317}{317}
\problemtext 在德布罗意与戴维孙--革末实验相关的四个论断中，
判断哪些真实反映微观粒子的波粒二象性。
\workings
光子具有波粒二象性，实物粒子也具有波粒二象性；
戴维孙--革末电子衍射实验验证了德布罗意物质波。
“电子只有粒子性”错误。故正确为命题1、2、3，选\textbf{B}。
\crosscheck 不可用经典轨道的粒子图像排除电子的干涉与衍射。

\sourceproblem{12.1.10}{317}{317}
\problemtext 电子由静止通过电势差$U$加速后德布罗意波长为$0.04$ nm，
求加速电压。
\workings
非相对论$p=\sqrt{2m_e eU}$，
$\lambda=h/p$；用常数式
$\lambda(\mathrm{nm})\simeq1.226/\sqrt{U(\mathrm V)}$，
\[
\boxed{U\simeq(1.226/0.04)^2\,\mathrm V
\simeq940\,\mathrm V}.
\]
选\textbf{D}。
\crosscheck 940 eV对511 keV的电子静能较小，非相对论近似合理。

\sourceproblem{12.1.11}{317}{317}
\problemtext 两种不同静质量的粒子具有相同德布罗意波长，
比较它们的动量、速度、动能与动能比。
\workings
由$\lambda=h/p$，两种粒子的\emph{动量大小}严格相同：
\[
\boxed{p_1=p_2=h/\lambda}.
\]
选\textbf{A}。
\crosscheck 同一动量下非相对论动能$p^2/(2m)$一般不同，
速度也一般不同。

\sourceproblem{12.1.12}{317}{317}
\problemtext 势阱$-a\le x\le a$内给归一化波函数
$\psi(x)=a^{-1/2}\cos(3\pi x/2a)$，求$x=5a/6$处概率密度。
\workings
$|\psi(x)|^2=a^{-1}\cos^2(3\pi x/2a)$，
代入$x=5a/6$，角度为$5\pi/4$，故
\[
\boxed{\rho(5a/6)=\frac1a\cos^2\frac{5\pi}{4}=\frac1{2a}}.
\]
选\textbf{A}。
\crosscheck 概率密度量纲是$\mathrm{length}^{-1}$，不是无量纲概率。

\sourceproblem{12.1.13}{317}{317}
\problemtext 判断不确定性原理是否禁止单独确定坐标、单独确定动量，
以及它是否适用于所有微观粒子。
\workings
不确定性原理限制的是同一态中位置与动量分量的不确定度\emph{同时}减小，
$\Delta x\,\Delta p_x\ge\hbar/2$；
单独某一量可以明确制备或测量。
因此正确的是命题(c)、(d)，选\textbf{B}。
\crosscheck 原题纸面把$\hbar$写成$h$样式；
采用标准量子力学常量$\hbar=h/(2\pi)$阐释，不将不同常量混用。

\sourceproblem{12.1.14}{317}{317}
\problemtext 电子处在原子K壳层（$n=1$），四种$(n,l,m_l,m_s)$组合中何者允许？
\workings
主量子数$n=1$时只能$l=0$，相应$m_l=0$，
自旋投影$m_s=\pm1/2$；
因此
\[
\boxed{(n,l,m_l,m_s)=(1,0,0,-1/2)}
\]
是选项\textbf{A}所列的允许态。
\crosscheck 源题选项B写$l=0,m_l=-1$，不满足$|m_l|\le l$；
C、D不属于$n=1$壳层。
